% !TeX program = xelatex
% !TeX encoding = UTF-8


\PassOptionsToPackage{dvipsnames}{xcolor}
% Add darkmode here to use the alternative palette.
\documentclass[11pt,a4paper,ragged2e,withhyper,rtl]{altacv}

\geometry{left=1.4cm,right=1.4cm,top=0.75cm,bottom=0.75cm,columnsep=0.85cm}
\usepackage{paracol}

\ifdarkmode
  \definecolor{PrimaryColor}{HTML}{EEE4D4}
  \definecolor{SecondaryColor}{HTML}{78C6BE}
  \definecolor{ThirdColor}{HTML}{E8B66B}
  \definecolor{BodyColor}{HTML}{D6D2CA}
  \definecolor{EmphasisColor}{HTML}{F5F0E8}
  \definecolor{BackgroundColor}{HTML}{18282B}
\else
  \definecolor{PrimaryColor}{HTML}{173F48}
  \definecolor{SecondaryColor}{HTML}{236C70}
  \definecolor{ThirdColor}{HTML}{C58A39}
  \definecolor{BodyColor}{HTML}{45494A}
  \definecolor{EmphasisColor}{HTML}{263A3E}
  \definecolor{BackgroundColor}{HTML}{FAF7F0}
\fi
\colorlet{name}{PrimaryColor}
\colorlet{tagline}{SecondaryColor}
\colorlet{heading}{PrimaryColor}
\colorlet{headingrule}{ThirdColor}
\colorlet{subheading}{SecondaryColor}
\colorlet{accent}{SecondaryColor}
\colorlet{emphasis}{EmphasisColor}
\colorlet{body}{BodyColor}
\pagecolor{BackgroundColor}

% Load xepersian LAST, after every other package, so bidi can patch them.
\usepackage{xepersian}
% \IfFontExistsTF{XB Yas}{%
%   \settextfont[Path = fonts/,
%     Extension = .ttf,
%     Scale=1.0]{XB Yas}
% }{%
%   \IfFontExistsTF{Amiri}{\settextfont[Scale=1.0]{Amiri}}{%
%     \IfFontExistsTF{Noto Naskh Arabic}{\settextfont[Scale=1.0]{Noto Naskh Arabic}}{%
%       \IfFontExistsTF{Tahoma}{\settextfont[Scale=1.0]{Tahoma}}{%
%         \ClassError{altacv}{No Persian/Arabic font found}{Install XB Yas, Amiri or Noto Naskh Arabic, then compile again.}
%       }
%     }
%   }
% }
\IfFileExists{fonts/XB Yas.ttf}{%
    \settextfont[
        Path = fonts/,
        Extension = .ttf,
        Scale = 1.0
    ]{XB Yas}%
}{%
    \IfFontExistsTF{Amiri}{%
        \settextfont[Scale=1.0]{Amiri}%
    }{%
        \IfFontExistsTF{Noto Naskh Arabic}{%
            \settextfont[Scale=1.0]{Noto Naskh Arabic}%
        }{%
            \IfFontExistsTF{Tahoma}{%
                \settextfont[Scale=1.0]{Tahoma}%
            }{%
                \ClassError{altacv}
                    {No Persian/Arabic font found}
                    {Upload XB Yas.ttf or install another Persian/Arabic font, then compile again.}%
            }%
        }%
    }%
}
\setlatintextfont{TeX Gyre Pagella}
\linespread{1.12}
\renewcommand{\locationname}{مکان}
\renewcommand{\datename}{زمان}
\hypersetup{pdftitle={رزومهٔ خیام نیشابوری},pdfauthor={عمر خیام نیشابوری}}

% Sizes live in the class's rtl defaults; override any of them here, e.g.:
% \renewcommand{\namefont}{\fontsize{26}{34}\selectfont\bfseries}
% Wrap inline English, URLs and Latin numbers in \cvlatin{...}.
% Persian and Arabic labels can be entered directly as UTF-8.

\begin{document}
\name{عُمَر خَیّام نیشابوری}
\tagline{ریاضی‌دان، ستاره‌شناس و شاعر؛ همه‌چیزدان}
\photoL{3.3cm}{khayyam}
\personalinfo{%
  غیاث‌الدین ابوالفتح عُمَر بن ابراهیم خَیّام نیشابوری\par
  \smallskip
  \location{نیشابور، ایران}
  \printinfo{\faClock}{قرن پنجم و ششم هجری}\par
  \smallskip
  % Fictional contact for the playful sample; linked fields stay LTR.
  \email{khayyam@example.com}
  \printinfo{\faBook}{رباعیات؛ نسخهٔ کاغذی لطفاً!}
}
\makecvheader

% Physical left column: sidebar (30%); right column: main CV (70%).
% Start writing in column 1 to put the main Persian content on the right.
% Both columns can still flow onto later pages when the CV grows.
\columnratio{0.30}
\begin{paracol}{2}
\switchcolumn[1]

\cvsection{دربارهٔ من}
\begin{quote}
روزها معادله حل می‌کنم، شب‌ها ستاره می‌شمارم و بین این دو، رباعی می‌گویم.
به پروژه‌هایی علاقه دارم که یا زمان را دقیق‌تر کنند یا زندگی را خوش‌تر؛
اگر هر دو شد، چه بهتر!
\end{quote}

\cvsection{تجربه‌ها}
\cvevent{پژوهشگرِ آسمان و تقویم}{همکاری در اصلاح تقویم جلالی}{دورهٔ سلجوقی}{اصفهان}
\begin{itemize}
  \item رصد و محاسبه برای هماهنگی تقویم با گردش فصل‌ها.
  \item کار گروهی با اهل نجوم؛ جلسه‌ها زیر آسمان برگزار می‌شد!
\end{itemize}
\divider

\cvevent{حل‌کنندهٔ مسائل سخت}{پژوهش در جبر و هندسه}{سال‌های پژوهش}{سمرقند و نیشابور}
\begin{itemize}
  \item بررسی معادلات درجهٔ سوم با کمک هندسه و مقاطع مخروطی.
  \item تبدیل پرسش‌های پیچیده به مسئله‌هایی که می‌شود درباره‌شان فکر کرد.
\end{itemize}

\cvsection{پروژه‌های منتخب}
\cvevent{تقویم جلالی}{وقتی نوروز باید سرِ وقت برسد}{}{ایران}
\begin{itemize}
  \item دقت در حساب روزها؛ بهار با حدس و گمان شروع نمی‌شود.
\end{itemize}
\divider

\cvevent{رسالهٔ جبر}{هندسه در خدمت معادله‌ها}{}{}
\begin{itemize}
  \item مطالعهٔ معادلات و جست‌وجوی راه‌حل‌های هندسی.
\end{itemize}
\divider

\cvevent{رباعیاتِ منسوب به خیام}{چهار مصراع، به اندازهٔ یک جهان}{}{}
\begin{itemize}
  \item تأمل در زمان، زندگی و لذتِ همین لحظه؛ با کمترین تعداد کلمات.
\end{itemize}

\cvsection{آموخته‌ها}
جبر، هندسه، نجوم و فلسفه؛ همراه با تمرینِ پرسیدنِ «چرا؟».
برای یاد گرفتن، هنوز هم یک مسئلهٔ خوب را به یک پاسخِ عجولانه ترجیح می‌دهم.

\switchcolumn[0]

\cvsection{جعبه‌ابزار}
% /true highlights a tag; the default is an ordinary tag.
\cvtags{جبر/true, هندسه/true, نجوم/true, فلسفه, شعر, تقویم, مقاطع مخروطی, رصد آسمان}

\cvsection{مهارت‌ها}
\cvskill{حل مسئله}{5}
\smallskip
\cvskill{Mathematics}{5}
\smallskip
\cvskill{نگاه به آسمان}{5}
\smallskip
\cvskill{کوتاه‌نویسی}{4.5}
\smallskip
\cvskill{جلسهٔ بی‌حاصل}{0.5}

\cvsection{زبان‌ها}
\cvlang{فارسی}{شعر و گفت‌وگو}\par
\smallskip
\cvlang{عربی}{نگارش علمی}\par
\smallskip
\cvlang{ریاضی}{بدون نیاز به مترجم}\par

\cvsection{افتخارها}
\cvachievement{\faCalendarCheck}{حسابِ روزها}{همکاری در تقویم جلالی}
\medskip
\cvachievement{\faFeather}{ماندگاریِ واژه‌ها}{رباعی‌هایی فراتر از زمان}

\cvsection{پس از کار}
تماشای آسمان، گپ با دوستان و پیدا کردنِ یک پرسش تازه برای فردا.

\cvsection{برای همکاری}
اگر مسئله‌تان دشوار است، خوشحال می‌شوم؛ اگر جلسه‌تان طولانی است، کمتر!

\bigskip
{\footnotesize\color{body!80}
این رزومه یک نمونهٔ طنز با الهام از خیام است؛ شرح شغل‌ها و راه ارتباطی خیالی‌اند.\par}

\end{paracol}
\end{document}
